% !TeX root = ../../Book2.tex
\section{热方程的结构}
\label{b2:sec:1203}

显式热核已经给出候选解。现在要分别检查正时间的导数、最大值控制和解的唯一性，并把光滑计算延伸到平方可积初值及时间可积的外力。每一种结论所用的解类都会保留在陈述中，尤其是全空间无穷远处的控制。

\subsection{光滑化效应}
\label{b2:subsec:120301}

热核的任意阶导数仍是多项式乘 Gauss 函数，因此正时间磨光不仅是定性结论。一次空间求导损失半个时间尺度，一次时间求导损失一个时间尺度；这是 \(x\sim\sqrt t\) 的抛物缩放直接产生的。

\Needspace{6\baselineskip}
\begin{theorem}[热流的导数估计]\label{b2:thm:12-smoothing}
设整数 \(d\geq1\)、\(\kappa>0\)，\(H_s(x)=(4\pi s)^{-d/2}e^{-|x|^2/(4s)}\)（\(s>0\)、\(x\in\mathbb R^d\)）。设 \(f\in L^p(\mathbb R^d)\)，\(1\leq p\leq q\leq\infty\)，\(\alpha\) 为多重指标，\(m\geq0\) 为整数。对 \(u(t)=H_{\kappa t}*f\)，有
\begin{equation}\label{b2:eq:12-smoothing}
 \|\partial_t^m\partial_x^\alpha u(t)\|_q
 \leq C t^{-m-|\alpha|/2-\frac d2(1/p-1/q)}\|f\|_p,
 \quad t>0.
\end{equation}
常数只依赖所列指数、维数与 \(\kappa\)，不依赖 \(t,f\)。
\end{theorem}
\begin{proof}
缩放与直接求导给出
\[
 \partial_t^m\partial_x^\alpha H_{\kappa t}(x)
 =t^{-m-|\alpha|/2-d/2}\Psi_{m,\alpha}(x/\sqrt t),
\]
其中 \(\Psi_{m,\alpha}\) 是快速衰减的光滑函数。因此其 \(L^r\) 范数为常数乘 \(t^{-m-|\alpha|/2-\frac d2(1-1/r)}\)。Young 不等式给出 \(p\to p\) 的导数界，逐点 Hölder 不等式给出 \(p\to\infty\) 的界。若 \(p<\infty\)、\(p\leq q<\infty\)，对导数函数 \(v\) 使用
\[
 \|v\|_q\leq\|v\|_p^{p/q}\cdot\|v\|_\infty^{1-p/q},
\]
由两个端点界即得所述指数。\(p=\infty\) 时只能取 \(q=\infty\)，已含在 Young 估计中。
\end{proof}

当 \(t\downarrow0\) 时，导数界可以发散，这是粗初值与正时间光滑性共存的方式。若初值已有某个 \(L^p\) 弱导数，可把该导数移到 \(f\) 上，就不必支付相应的时间奇性。取一维跳跃初温 \(f=\mathbf1_{(0,\infty)}\)，解为
\[
 u(t,x)=\dfrac1{\sqrt\pi}\int_{-\infty}^{x/\sqrt{4\kappa t}}e^{-s^2}\,\mathrm{d}s,
 \quad u_x(t,x)=H_{\kappa t}(x).
\]
跃变被摊开到 \(\sqrt{\kappa t}\) 的过渡层，梯度峰值为 \((4\pi\kappa t)^{-1/2}\)。这同时给出不能在整个 \(L^\infty\) 中要求初值一致收敛的具体例子。\Cref{b2:fig:12-heat-step}将三个正时间的剖面与初温并列：曲线随时间变缓，过渡层随之变宽；在初温为零的一侧，正时间已处处出现非零的细小尾部。
% Step heat profiles: values sampled directly from the proved Gaussian integral, kappa=1.
\begin{figure}[htbp]
\centering
\begin{tikzpicture}[x=1.1cm,y=3.4cm]
\draw[->,AnalysisGray] (-3.35,0)--(3.35,0) node[right]{\(x\)};
\draw[->,AnalysisGray] (-3.3,0)--(-3.3,1.15) node[above]{\(u\)};
\foreach \x in {-3,-2,-1,0,1,2,3}{
 \draw (\x,0)--(\x,-.02) node[below,font=\small]{\(\x\)};
}
\foreach \y/\lab in {.25/{\frac14},.5/{\frac12},.75/{\frac34},1/1}{
 \draw (-3.3,\y)--(-3.36,\y) node[left,font=\small]{\(\lab\)};
}
\draw[AnalysisGray!35,thin,densely dotted] (-3.28,.5)--(3.15,.5);
\draw[black,thick,densely dotted] (-3,0)--(0,0);
\draw[black,thick,densely dotted] (0,1)--(3,1);
\fill (0,0) circle(1.3pt);
\draw[fill=white] (0,1) circle(1.3pt);
\draw[AnalysisBlue,very thick] plot coordinates {
 (-3.0000,0.00000000) (-2.9700,0.00000000) (-2.9400,0.00000000) (-2.9100,0.00000000) (-2.8800,0.00000000) (-2.8500,0.00000000) (-2.8200,0.00000000) (-2.7900,0.00000000)
 (-2.7600,0.00000000) (-2.7300,0.00000000) (-2.7000,0.00000000) (-2.6700,0.00000000) (-2.6400,0.00000000) (-2.6100,0.00000000) (-2.5800,0.00000000) (-2.5500,0.00000000)
 (-2.5200,0.00000000) (-2.4900,0.00000000) (-2.4600,0.00000000) (-2.4300,0.00000000) (-2.4000,0.00000000) (-2.3700,0.00000000) (-2.3400,0.00000000) (-2.3100,0.00000000)
 (-2.2800,0.00000000) (-2.2500,0.00000000) (-2.2200,0.00000000) (-2.1900,0.00000000) (-2.1600,0.00000000) (-2.1300,0.00000000) (-2.1000,0.00000000) (-2.0700,0.00000000)
 (-2.0400,0.00000000) (-2.0100,0.00000001) (-1.9800,0.00000001) (-1.9500,0.00000002) (-1.9200,0.00000003) (-1.8900,0.00000005) (-1.8600,0.00000007) (-1.8300,0.00000011)
 (-1.8000,0.00000018) (-1.7700,0.00000028) (-1.7400,0.00000043) (-1.7100,0.00000066) (-1.6800,0.00000101) (-1.6500,0.00000153) (-1.6200,0.00000230) (-1.5900,0.00000344)
 (-1.5600,0.00000511) (-1.5300,0.00000754) (-1.5000,0.00001105) (-1.4700,0.00001607) (-1.4400,0.00002321) (-1.4100,0.00003330) (-1.3800,0.00004746) (-1.3500,0.00006717)
 (-1.3200,0.00009441) (-1.2900,0.00013180) (-1.2600,0.00018275) (-1.2300,0.00025168) (-1.2000,0.00034426) (-1.1700,0.00046772) (-1.1400,0.00063117) (-1.1100,0.00084603)
 (-1.0800,0.00112642) (-1.0500,0.00148973) (-1.0200,0.00195710) (-0.9900,0.00255400) (-0.9600,0.00331089) (-0.9300,0.00426376) (-0.9000,0.00545475) (-0.8700,0.00693270)
 (-0.8400,0.00875361) (-0.8100,0.01098096) (-0.7800,0.01368594) (-0.7500,0.01694743) (-0.7200,0.02085172) (-0.6900,0.02549198) (-0.6600,0.03096742) (-0.6300,0.03738203)
 (-0.6000,0.04484301) (-0.5700,0.05345884) (-0.5400,0.06333692) (-0.5100,0.07458099) (-0.4800,0.08728818) (-0.4500,0.10154589) (-0.4200,0.11742864) (-0.3900,0.13499478)
 (-0.3600,0.15428344) (-0.3300,0.17531166) (-0.3000,0.19807195) (-0.2700,0.22253037) (-0.2400,0.24862516) (-0.2100,0.27626619) (-0.1800,0.30533515) (-0.1500,0.33568662)
 (-0.1200,0.36714997) (-0.0900,0.39953208) (-0.0600,0.43262082) (-0.0300,0.46618920) (0.0000,0.50000000) (0.0300,0.53381080) (0.0600,0.56737918) (0.0900,0.60046792)
 (0.1200,0.63285003) (0.1500,0.66431338) (0.1800,0.69466485) (0.2100,0.72373381) (0.2400,0.75137484) (0.2700,0.77746963) (0.3000,0.80192805) (0.3300,0.82468834)
 (0.3600,0.84571656) (0.3900,0.86500522) (0.4200,0.88257136) (0.4500,0.89845411) (0.4800,0.91271182) (0.5100,0.92541901) (0.5400,0.93666308) (0.5700,0.94654116)
 (0.6000,0.95515699) (0.6300,0.96261797) (0.6600,0.96903258) (0.6900,0.97450802) (0.7200,0.97914828) (0.7500,0.98305257) (0.7800,0.98631406) (0.8100,0.98901904)
 (0.8400,0.99124639) (0.8700,0.99306730) (0.9000,0.99454525) (0.9300,0.99573624) (0.9600,0.99668911) (0.9900,0.99744600) (1.0200,0.99804290) (1.0500,0.99851027)
 (1.0800,0.99887358) (1.1100,0.99915397) (1.1400,0.99936883) (1.1700,0.99953228) (1.2000,0.99965574) (1.2300,0.99974832) (1.2600,0.99981725) (1.2900,0.99986820)
 (1.3200,0.99990559) (1.3500,0.99993283) (1.3800,0.99995254) (1.4100,0.99996670) (1.4400,0.99997679) (1.4700,0.99998393) (1.5000,0.99998895) (1.5300,0.99999246)
 (1.5600,0.99999489) (1.5900,0.99999656) (1.6200,0.99999770) (1.6500,0.99999847) (1.6800,0.99999899) (1.7100,0.99999934) (1.7400,0.99999957) (1.7700,0.99999972)
 (1.8000,0.99999982) (1.8300,0.99999989) (1.8600,0.99999993) (1.8900,0.99999995) (1.9200,0.99999997) (1.9500,0.99999998) (1.9800,0.99999999) (2.0100,0.99999999)
 (2.0400,1.00000000) (2.0700,1.00000000) (2.1000,1.00000000) (2.1300,1.00000000) (2.1600,1.00000000) (2.1900,1.00000000) (2.2200,1.00000000) (2.2500,1.00000000)
 (2.2800,1.00000000) (2.3100,1.00000000) (2.3400,1.00000000) (2.3700,1.00000000) (2.4000,1.00000000) (2.4300,1.00000000) (2.4600,1.00000000) (2.4900,1.00000000)
 (2.5200,1.00000000) (2.5500,1.00000000) (2.5800,1.00000000) (2.6100,1.00000000) (2.6400,1.00000000) (2.6700,1.00000000) (2.7000,1.00000000) (2.7300,1.00000000)
 (2.7600,1.00000000) (2.7900,1.00000000) (2.8200,1.00000000) (2.8500,1.00000000) (2.8800,1.00000000) (2.9100,1.00000000) (2.9400,1.00000000) (2.9700,1.00000000)
 (3.0000,1.00000000)
};
\draw[AnalysisGray,thick,densely dashed] plot coordinates {
 (-3.0000,0.00001105) (-2.9700,0.00001333) (-2.9400,0.00001607) (-2.9100,0.00001933) (-2.8800,0.00002321) (-2.8500,0.00002783) (-2.8200,0.00003330) (-2.7900,0.00003979)
 (-2.7600,0.00004746) (-2.7300,0.00005651) (-2.7000,0.00006717) (-2.6700,0.00007970) (-2.6400,0.00009441) (-2.6100,0.00011164) (-2.5800,0.00013180) (-2.5500,0.00015533)
 (-2.5200,0.00018275) (-2.4900,0.00021464) (-2.4600,0.00025168) (-2.4300,0.00029460) (-2.4000,0.00034426) (-2.3700,0.00040161) (-2.3400,0.00046772) (-2.3100,0.00054379)
 (-2.2800,0.00063117) (-2.2500,0.00073136) (-2.2200,0.00084603) (-2.1900,0.00097703) (-2.1600,0.00112642) (-2.1300,0.00129649) (-2.1000,0.00148973) (-2.0700,0.00170892)
 (-2.0400,0.00195710) (-2.0100,0.00223758) (-1.9800,0.00255400) (-1.9500,0.00291033) (-1.9200,0.00331089) (-1.8900,0.00376034) (-1.8600,0.00426376) (-1.8300,0.00482660)
 (-1.8000,0.00545475) (-1.7700,0.00615453) (-1.7400,0.00693270) (-1.7100,0.00779650) (-1.6800,0.00875361) (-1.6500,0.00981221) (-1.6200,0.01098096) (-1.5900,0.01226899)
 (-1.5600,0.01368594) (-1.5300,0.01524190) (-1.5000,0.01694743) (-1.4700,0.01881355) (-1.4400,0.02085172) (-1.4100,0.02307378) (-1.3800,0.02549198) (-1.3500,0.02811890)
 (-1.3200,0.03096742) (-1.2900,0.03405068) (-1.2600,0.03738203) (-1.2300,0.04097495) (-1.2000,0.04484301) (-1.1700,0.04899980) (-1.1400,0.05345884) (-1.1100,0.05823349)
 (-1.0800,0.06333692) (-1.0500,0.06878195) (-1.0200,0.07458099) (-0.9900,0.08074597) (-0.9600,0.08728818) (-0.9300,0.09421822) (-0.9000,0.10154589) (-0.8700,0.10928008)
 (-0.8400,0.11742864) (-0.8100,0.12599836) (-0.7800,0.13499478) (-0.7500,0.14442218) (-0.7200,0.15428344) (-0.6900,0.16457997) (-0.6600,0.17531166) (-0.6300,0.18647678)
 (-0.6000,0.19807195) (-0.5700,0.21009210) (-0.5400,0.22253037) (-0.5100,0.23537819) (-0.4800,0.24862516) (-0.4500,0.26225914) (-0.4200,0.27626619) (-0.3900,0.29063065)
 (-0.3600,0.30533515) (-0.3300,0.32036067) (-0.3000,0.33568662) (-0.2700,0.35129089) (-0.2400,0.36714997) (-0.2100,0.38323904) (-0.1800,0.39953208) (-0.1500,0.41600201)
 (-0.1200,0.43262082) (-0.0900,0.44935970) (-0.0600,0.46618920) (-0.0300,0.48307939) (0.0000,0.50000000) (0.0300,0.51692061) (0.0600,0.53381080) (0.0900,0.55064030)
 (0.1200,0.56737918) (0.1500,0.58399799) (0.1800,0.60046792) (0.2100,0.61676096) (0.2400,0.63285003) (0.2700,0.64870911) (0.3000,0.66431338) (0.3300,0.67963933)
 (0.3600,0.69466485) (0.3900,0.70936935) (0.4200,0.72373381) (0.4500,0.73774086) (0.4800,0.75137484) (0.5100,0.76462181) (0.5400,0.77746963) (0.5700,0.78990790)
 (0.6000,0.80192805) (0.6300,0.81352322) (0.6600,0.82468834) (0.6900,0.83542003) (0.7200,0.84571656) (0.7500,0.85557782) (0.7800,0.86500522) (0.8100,0.87400164)
 (0.8400,0.88257136) (0.8700,0.89071992) (0.9000,0.89845411) (0.9300,0.90578178) (0.9600,0.91271182) (0.9900,0.91925403) (1.0200,0.92541901) (1.0500,0.93121805)
 (1.0800,0.93666308) (1.1100,0.94176651) (1.1400,0.94654116) (1.1700,0.95100020) (1.2000,0.95515699) (1.2300,0.95902505) (1.2600,0.96261797) (1.2900,0.96594932)
 (1.3200,0.96903258) (1.3500,0.97188110) (1.3800,0.97450802) (1.4100,0.97692622) (1.4400,0.97914828) (1.4700,0.98118645) (1.5000,0.98305257) (1.5300,0.98475810)
 (1.5600,0.98631406) (1.5900,0.98773101) (1.6200,0.98901904) (1.6500,0.99018779) (1.6800,0.99124639) (1.7100,0.99220350) (1.7400,0.99306730) (1.7700,0.99384547)
 (1.8000,0.99454525) (1.8300,0.99517340) (1.8600,0.99573624) (1.8900,0.99623966) (1.9200,0.99668911) (1.9500,0.99708967) (1.9800,0.99744600) (2.0100,0.99776242)
 (2.0400,0.99804290) (2.0700,0.99829108) (2.1000,0.99851027) (2.1300,0.99870351) (2.1600,0.99887358) (2.1900,0.99902297) (2.2200,0.99915397) (2.2500,0.99926864)
 (2.2800,0.99936883) (2.3100,0.99945621) (2.3400,0.99953228) (2.3700,0.99959839) (2.4000,0.99965574) (2.4300,0.99970540) (2.4600,0.99974832) (2.4900,0.99978536)
 (2.5200,0.99981725) (2.5500,0.99984467) (2.5800,0.99986820) (2.6100,0.99988836) (2.6400,0.99990559) (2.6700,0.99992030) (2.7000,0.99993283) (2.7300,0.99994349)
 (2.7600,0.99995254) (2.7900,0.99996021) (2.8200,0.99996670) (2.8500,0.99997217) (2.8800,0.99997679) (2.9100,0.99998067) (2.9400,0.99998393) (2.9700,0.99998667)
 (3.0000,0.99998895)
};
\draw[black!55,thick,dash dot] plot coordinates {
 (-3.0000,0.01694743) (-2.9700,0.01785971) (-2.9400,0.01881355) (-2.9100,0.01981040) (-2.8800,0.02085172) (-2.8500,0.02193900) (-2.8200,0.02307378) (-2.7900,0.02425759)
 (-2.7600,0.02549198) (-2.7300,0.02677855) (-2.7000,0.02811890) (-2.6700,0.02951465) (-2.6400,0.03096742) (-2.6100,0.03247888) (-2.5800,0.03405068) (-2.5500,0.03568451)
 (-2.5200,0.03738203) (-2.4900,0.03914494) (-2.4600,0.04097495) (-2.4300,0.04287374) (-2.4000,0.04484301) (-2.3700,0.04688447) (-2.3400,0.04899980) (-2.3100,0.05119070)
 (-2.2800,0.05345884) (-2.2500,0.05580588) (-2.2200,0.05823349) (-2.1900,0.06074330) (-2.1600,0.06333692) (-2.1300,0.06601595) (-2.1000,0.06878195) (-2.0700,0.07163646)
 (-2.0400,0.07458099) (-2.0100,0.07761701) (-1.9800,0.08074597) (-1.9500,0.08396924) (-1.9200,0.08728818) (-1.8900,0.09070409) (-1.8600,0.09421822) (-1.8300,0.09783178)
 (-1.8000,0.10154589) (-1.7700,0.10536165) (-1.7400,0.10928008) (-1.7100,0.11330211) (-1.6800,0.11742864) (-1.6500,0.12166048) (-1.6200,0.12599836) (-1.5900,0.13044294)
 (-1.5600,0.13499478) (-1.5300,0.13965440) (-1.5000,0.14442218) (-1.4700,0.14929845) (-1.4400,0.15428344) (-1.4100,0.15937726) (-1.3800,0.16457997) (-1.3500,0.16989149)
 (-1.3200,0.17531166) (-1.2900,0.18084021) (-1.2600,0.18647678) (-1.2300,0.19222089) (-1.2000,0.19807195) (-1.1700,0.20402929) (-1.1400,0.21009210) (-1.1100,0.21625946)
 (-1.0800,0.22253037) (-1.0500,0.22890370) (-1.0200,0.23537819) (-0.9900,0.24195250) (-0.9600,0.24862516) (-0.9300,0.25539461) (-0.9000,0.26225914) (-0.8700,0.26921697)
 (-0.8400,0.27626619) (-0.8100,0.28340479) (-0.7800,0.29063065) (-0.7500,0.29794155) (-0.7200,0.30533515) (-0.6900,0.31280903) (-0.6600,0.32036067) (-0.6300,0.32798744)
 (-0.6000,0.33568662) (-0.5700,0.34345540) (-0.5400,0.35129089) (-0.5100,0.35919010) (-0.4800,0.36714997) (-0.4500,0.37516736) (-0.4200,0.38323904) (-0.3900,0.39136173)
 (-0.3600,0.39953208) (-0.3300,0.40774667) (-0.3000,0.41600201) (-0.2700,0.42429459) (-0.2400,0.43262082) (-0.2100,0.44097708) (-0.1800,0.44935970) (-0.1500,0.45776499)
 (-0.1200,0.46618920) (-0.0900,0.47462860) (-0.0600,0.48307939) (-0.0300,0.49153779) (0.0000,0.50000000) (0.0300,0.50846221) (0.0600,0.51692061) (0.0900,0.52537140)
 (0.1200,0.53381080) (0.1500,0.54223501) (0.1800,0.55064030) (0.2100,0.55902292) (0.2400,0.56737918) (0.2700,0.57570541) (0.3000,0.58399799) (0.3300,0.59225333)
 (0.3600,0.60046792) (0.3900,0.60863827) (0.4200,0.61676096) (0.4500,0.62483264) (0.4800,0.63285003) (0.5100,0.64080990) (0.5400,0.64870911) (0.5700,0.65654460)
 (0.6000,0.66431338) (0.6300,0.67201256) (0.6600,0.67963933) (0.6900,0.68719097) (0.7200,0.69466485) (0.7500,0.70205845) (0.7800,0.70936935) (0.8100,0.71659521)
 (0.8400,0.72373381) (0.8700,0.73078303) (0.9000,0.73774086) (0.9300,0.74460539) (0.9600,0.75137484) (0.9900,0.75804750) (1.0200,0.76462181) (1.0500,0.77109630)
 (1.0800,0.77746963) (1.1100,0.78374054) (1.1400,0.78990790) (1.1700,0.79597071) (1.2000,0.80192805) (1.2300,0.80777911) (1.2600,0.81352322) (1.2900,0.81915979)
 (1.3200,0.82468834) (1.3500,0.83010851) (1.3800,0.83542003) (1.4100,0.84062274) (1.4400,0.84571656) (1.4700,0.85070155) (1.5000,0.85557782) (1.5300,0.86034560)
 (1.5600,0.86500522) (1.5900,0.86955706) (1.6200,0.87400164) (1.6500,0.87833952) (1.6800,0.88257136) (1.7100,0.88669789) (1.7400,0.89071992) (1.7700,0.89463835)
 (1.8000,0.89845411) (1.8300,0.90216822) (1.8600,0.90578178) (1.8900,0.90929591) (1.9200,0.91271182) (1.9500,0.91603076) (1.9800,0.91925403) (2.0100,0.92238299)
 (2.0400,0.92541901) (2.0700,0.92836354) (2.1000,0.93121805) (2.1300,0.93398405) (2.1600,0.93666308) (2.1900,0.93925670) (2.2200,0.94176651) (2.2500,0.94419412)
 (2.2800,0.94654116) (2.3100,0.94880930) (2.3400,0.95100020) (2.3700,0.95311553) (2.4000,0.95515699) (2.4300,0.95712626) (2.4600,0.95902505) (2.4900,0.96085506)
 (2.5200,0.96261797) (2.5500,0.96431549) (2.5800,0.96594932) (2.6100,0.96752112) (2.6400,0.96903258) (2.6700,0.97048535) (2.7000,0.97188110) (2.7300,0.97322145)
 (2.7600,0.97450802) (2.7900,0.97574241) (2.8200,0.97692622) (2.8500,0.97806100) (2.8800,0.97914828) (2.9100,0.98018960) (2.9400,0.98118645) (2.9700,0.98214029)
 (3.0000,0.98305257)
};
\draw[AnalysisGray,thin,densely dashed] (-0.95387255,-.24)--(-0.95387255,-.16)
 (-0.95387255,.02)--(-0.95387255,.25);
\draw[AnalysisGray,thin,densely dashed] (0.95387255,-.24)--(0.95387255,-.16)
 (0.95387255,.02)--(0.95387255,.75);
\fill[black!55] (-0.95387255,.25) circle(1.2pt);
\fill[black!55] (0.95387255,.75) circle(1.2pt);
\draw[<->,AnalysisGray] (-0.95387255,-.24)--(0.95387255,-.24);
\node[align=center,font=\small] at (0,-.39){固定温度比例间的过渡宽度\(\asymp\sqrt{\kappa t}\)};
\draw[black,thick,densely dotted] (3.9,1.02)--(4.5,1.02);
\node[anchor=west,font=\small] at (4.6,1.02){\(t=0\)};
\draw[AnalysisBlue,very thick] (3.9,.82)--(4.5,.82);
\node[anchor=west,font=\small] at (4.6,.82){\(t=1/16\)};
\draw[AnalysisGray,thick,densely dashed] (3.9,.62)--(4.5,.62);
\node[anchor=west,font=\small] at (4.6,.62){\(t=1/4\)};
\draw[black!55,thick,dash dot] (3.9,.42)--(4.5,.42);
\node[anchor=west,font=\small] at (4.6,.42){\(t=1\)};
\end{tikzpicture}
\caption[阶跃初温的光滑化与扩散尺度]{阶跃初温的光滑化与扩散尺度。图取\(\kappa=1\)，点线为\(f=\mathbf1_{(0,\infty)}\)，其余三条曲线按正文的 Gaussian 积分计算；下方标出\(t=1\)时温度为\(1/4\)、\(3/4\)的两点间距。任意固定温度比例间的过渡宽度与\(\sqrt{\kappa t}\)成正比。每个正时间剖面均光滑，且对每个有限\(x\)都有\(0<u(t,x)<1\)：原先零温的任意远处已出现正温度，但远离界面的影响可以极弱，图上线条与水平线近乎重合不表示尾部为零。}
\label{b2:fig:12-heat-step}
\end{figure}


\subsection{最大值原理}
\label{b2:subsec:120302}

最大值原理从微分方程本身控制解，不依赖可显式计算的核。本小节讨论实值函数。若 \(\Omega\subset\mathbb R^d\) 是有界开集、\(T>0\)，记 \(Q_T=(0,T]\times\Omega\)，其\term[paowu-bianjie]{抛物边界}[parabolic boundary]为
\[
 \partial_pQ_T=(\{0\}\times\overline\Omega)
 \cup([0,T]\times\partial\Omega).
\]
上端面 \(t=T\) 的内部点不属于这条边界：它们仍受到先前时间的数据控制。

\Needspace{6\baselineskip}
\begin{theorem}[抛物最大值原理]\label{b2:thm:12-maximum}
设整数 \(d\geq1\)、\(\kappa,T>0\)，\(\Omega\subset\mathbb R^d\) 为非空有界开集，\(\partial_pQ_T=(\{0\}\times\overline\Omega)\cup([0,T]\times\partial\Omega)\)。设实值 \(u\in C([0,T]\times\overline\Omega)\)，在正时间内有一阶时间导数和二阶空间导数，且 \(u_t-\kappa\Delta u\leq0\)。则
\[
 \max_{[0,T]\times\overline\Omega}u
 \leq\max_{\partial_pQ_T}u.
\]
导数在 \(t=T\) 取左侧时间导数即可。
\end{theorem}
\begin{proof}
固定 \(\varepsilon>0\)，令 \(v=u-\varepsilon t\)，则 \(v_t-\kappa\Delta v<0\)。若 \(v\) 的最大值严格超过抛物边界的最大值，紧性使它在某个正时间、空间内部的点取得最大值。该点有 \(\nabla_xv=0\)、\(\Delta v\leq0\)，且相对于过去时间有 \(v_t\geq0\)，故 \(v_t-\kappa\Delta v\geq0\)，矛盾。于是 \(u(t,x)\leq\max_{\partial_pQ_T}u+\varepsilon T\)，令 \(\varepsilon\downarrow0\)。
\end{proof}

这称为\term[paowu-zuidazhi-yuanli]{抛物最大值原理}[parabolic maximum principle]。对 \(-u\) 使用它得到最小值控制，对两个解之差使用它得到比较原理。特别是同样的 Dirichlet 边界条件下，初值误差的一致范数不会增长。有非齐次源项时，若 \(|F|\leq M\)，比较函数可加上 \(Mt\)，从而得到初边值最大误差再加 \(MT\) 的估计。

\ProseParagraphBreak

无界空间不能直接援引紧性取得最大值，也不能忽略远处边界。对每个有限时间带上一致有界的经典次解，可以用障碍函数
\[
 \Phi(t,x)=1+|x|^2+2d\kappa t,
 \quad\Phi_t-\kappa\Delta\Phi=0
\]
补上这一步。若初值 \(u(0,x)\leq0\)，在半径 \(R\) 的球柱中比较 \(u\) 与 \(\varepsilon\Phi\)。有界性使足够大的侧边界满足 \(u\leq\varepsilon\Phi\)，初始面也满足；有界区域最大值原理于是给出任意固定点的同一不等式。先令 \(R\) 足够大，再令 \(\varepsilon\downarrow0\)，得到 \(u\leq0\)。

\subsection{唯一性}
\label{b2:subsec:120303}

上面的障碍证明给出了全空间有界经典解的唯一性：在每个 \([0,T]\times\mathbb R^d\) 上有共同有限界、连续取得相同初值的两个解必相等。该结论不覆盖任意增长的经典解。空间无穷远的限制参与证明，不能从有界区域的唯一性直接删去。

\ProseParagraphBreak

对平方可积初值，另一种自然解类由能量决定。以下版本只要求分布方程与时间上的 \(L^2\) 连续性，因此也覆盖此前构造的粗初值解。

\Needspace{6\baselineskip}
\begin{theorem}[平方可积解的唯一性]\label{b2:thm:12-l2-uniqueness}
设整数 \(d\geq1\)、\(\kappa>0\)，\(H_s(x)=(4\pi s)^{-d/2}e^{-|x|^2/(4s)}\)（\(s>0\)、\(x\in\mathbb R^d\)）。设 \(T>0\)、\(f\in L^2(\mathbb R^d;\mathbb C)\)。若 \(u\in C([0,T];L^2(\mathbb R^d))\) 在分布意义下满足 \(u_t=\kappa\Delta u\)，且 \(u(0)=f\)，则 \(u(t)=H_{\kappa t}*f\)。此外，核解满足
\begin{equation}\label{b2:eq:12-whole-energy}
 \|u(T)\|_2^2+2\kappa\int_0^T\|\nabla u(t)\|_2^2\,\mathrm{d}t
 =\|f\|_2^2.
\end{equation}
\end{theorem}
\begin{proof}
空间 Fourier 变换在 \(L^2\) 上酉。先说明可以对方程使用 Schwartz 空间检验函数：将它乘以趋于一的空间紧支集截断，函数及二阶导数都在 \(L^2\) 中收敛；\(u\) 在紧时间区间上范数有界，所以分布恒等式传到极限。因此空间变换后有 \(\partial_t\widehat u+a\widehat u=0\)，其中 \(a(\xi)=4\pi^2\kappa|\xi|^2\)。

取任意紧支集光滑频率函数 \(\chi\)。在其支集上，\(e^{ta(\xi)}\) 及所需导数有界。乘以这个因子得到
\[
 \partial_t\left(e^{ta}\chi\widehat u(t)\right)=0.
\]
对每个光滑紧支集的频率测试函数配对，所得标量分布导数为零，故是常数；它又随时间连续，常数由初值确定。因此 \(\chi\widehat u(t)=e^{-ta}\chi\widehat f\)。取覆盖所有频率的截断列，得到 \(\widehat u(t)=e^{-ta}\widehat f\) 于 \(L^2\)，从而唯一。

最后对每个频率有
\[
 e^{-2Ta}|\widehat f|^2
 +2\int_0^Tae^{-2ta}|\widehat f|^2\,\mathrm{d}t
 =|\widehat f|^2.
\]
非负积分可换序；Plancherel 与平方可积弱导数的频率刻画把第二项变为 \(2\kappa\int\|\nabla u\|_2^2\)，得到能量恒等式。
\end{proof}

有界经典类与 \(C_tL^2_x\) 类各有独立的远处控制：前者允许非零常数初温，后者容许无界但平方可积的初温。二者有交集，却不相互包含。对一个给定初值问题，先明确解类，再谈唯一性。

\Needspace{6\baselineskip}
\begin{proposition}[粗外力与热能量]\label{b2:prop:12-forced-energy}
设整数 \(d\geq1\)、\(\kappa>0\)，\(H_s(x)=(4\pi s)^{-d/2}e^{-|x|^2/(4s)}\)（\(s>0\)、\(x\in\mathbb R^d\)）。设 \(T>0\)、\(f\in L^2(\mathbb R^d;\mathbb C)\)、\(F\in L^1((0,T);L^2(\mathbb R^d))\)。式\eqref{b2:eq:12-duhamel}的积分以 \(L^2\) 值积分解释，给出唯一的 \(C([0,T];L^2)\) 分布解，满足 \(u_t-\kappa\Delta u=F\) 及初值 \(f\)。它满足
\[
 \|u(t)\|_2\leq\|f\|_2+\int_0^t\|F(s)\|_2\,\mathrm{d}s,
 \quad \nabla u\in L^2((0,T)\times\mathbb R^d),
\]
以及能量恒等式
\[
 \|u(t)\|_2^2+2\kappa\int_0^t\|\nabla u(s)\|_2^2\,\mathrm{d}s
 =\|f\|_2^2+2\operatorname{Re}\int_0^t\langle F(s),u(s)\rangle_{L^2}\,\mathrm{d}s.
\]
\end{proposition}
\begin{proof}
热卷积在 \(L^2\) 中为收缩，并在零时强趋于恒等。先对简单外力定义积分，再以积分范数逼近，得到公式及第一项估计。时间连续性可先对有限个空间取值的简单函数检验：各段积分由核的强连续性和控制收敛得到连续，端点处新增积分的范数趋于零。一般外力的误差在整个 \([0,T]\) 上不超过其 \(L^1_tL^2_x\) 范数，故同样连续，且积分项在零时趋于零。

选光滑紧支集 \(f_j\to f\) 于 \(L^2\)，并选光滑时空外力 \(F_j\to F\) 于 \(L^1_tL^2_x\)，使每个固定 \(j\) 的 \(F_j\) 在整个时间带上有紧空间支集，紧集允许随 \(j\) 扩大。后者可先由简单函数逼近，再分别磨光有限个空间取值及时间示性函数。对应光滑解 \(u_j\) 在 \(C_tL^2_x\) 中趋于 \(u\)，方程的测试等式随之传到极限。对这些光滑解乘以 \(\overline{u_j}\)、取实部并作空间积分，得到所称能量恒等式。全空间分部积分可先乘标准截断 \(\chi_R^2\)：额外项的模至多为 \(CR^{-1}\|u_j\|_2\cdot\|\nabla u_j\|_2\)。对每个固定 \(j\)，初值与外力的光滑紧支集逼近使 \(\|u_j(t)\|_2\)、\(\|\nabla u_j(t)\|_2\) 在该有限时间带上有界，这也可将一次空间导数移到平滑初值和外力上后用 Young 不等式得到。因此截断通量的时间积分趋于零，再令 \(R\to\infty\) 即得上述恒等式；这里不要求对所有 \(j\) 共用一个空间支集。

还需证明梯度收敛，不能仅由零阶收敛推断。记两近似解之差为 \(v\)，初值差为 \(h\)，外力差为 \(G\)。差解的能量恒等式给出
\[
 2\kappa\int_0^T\|\nabla v(s)\|_2^2\,\mathrm{d}s
 \leq\|h\|_2^2+2\sup_{s\leq T}\|v(s)\|_2
                    \int_0^T\|G(s)\|_2\,\mathrm{d}s.
\]
第一项估计又控制 \(\sup\|v\|_2\)，所以梯度在时空 \(L^2\) 中为 Cauchy 列。对空间测试函数分部积分取极限，将其极限识别为 \(\nabla u\)。于是能量的耗散积分收敛；外力功由外力的积分范数收敛及 \(u_j\) 的一致 \(L^2\) 收敛传到极限。所得等式对每个 \(t\in[0,T]\) 成立。两个同数据解之差满足齐次方程及零初值，\cref{b2:thm:12-l2-uniqueness}给出唯一性。
\end{proof}

\subsection{保正性}
\label{b2:subsec:120304}

设 \(\kappa>0\)、\(1\le p\le\infty\)、\(f\in L^p(\mathbb R^d;\mathbb R)\)。热核严格为正。因此若 \(f\geq0\) 几乎处处，则 \(H_{\kappa t}*f\geq0\)；若 \(f\) 不是零等价类，某个正测度集合上有 \(f>0\)，对每个 \(t>0,x\) 的积分都严格为正。这一性质称为\term[baozheng-xing]{保正性}[positivity preservation]。即使初温支集很小，任何正时间全空间每一点都受到影响；热传播没有波方程的有限传播速度。

\ProseParagraphBreak

在有限区间 \((0,\ell)\)、\(\ell,\kappa>0\) 上，纯热方程 \(u_t=\kappa u_{xx}\) 的非负实初温 \(f\in L^2(0,\ell;\mathbb R)\) 也产生非负解。对零 Dirichlet 条件，先用非负光滑紧支集函数逼近初值，最大值原理给出各近似解非负；正时间核的 \(L^2\to L^\infty\) 界让它们一致趋于目标解，故极限非负。零 Neumann 的镜像核 \(K_N\) 本身为正。Robin 非负系数时，可以直接使用负部能量，沿用\secref{b2:sec:1201}的通量约定，边界写为 \(-\kappa u_x(0)+h_0u(0)=\kappa u_x(\ell)+h_\ell u(\ell)=0\)，\(h_0,h_\ell\geq0\)。正时间有 \(u(t)\in D(A)\)、\(u_t(t)\in H\)，因此不需预设初温有经典导数。

\ProseParagraphBreak

令 \(u_-=\max\{-u,0\}\)。绝对连续函数与这个 Lipschitz 函数复合后仍绝对连续，且 \((u_-)'=-u'\cdot\mathbf1_{\{u<0\}}\) 几乎处处；在零水平集上 \(u'=0\) 几乎处处，这可在可微且为该水平集密度点的位置沿零点差商取得。故 \(u_-\in V\)。对标量函数 \(J(z)=z_-^2/2\)，导数为 \(-z_-\)，且导数为 Lipschitz 函数，因此 Taylor 余项被 \(|h|^2/2\) 控制。积分这个余项，再使用 \(H\) 中的时间可微性，得到 \(\dfrac{\mathrm{d}}{\mathrm{d}t}\int J(u)=-\langle u_t,u_-\rangle\)。将方程与 \(u_-\) 配对，并使用 \(\mathfrak a(u,u_-)=-\kappa\|(u_-)'\|_2^2-\sum_s h_s|u_-(s)|^2\)，便有
\[
 \dfrac12\dfrac{\mathrm{d}}{\mathrm{d}t}\|u_-\|_2^2
 =-\kappa\|(u_-)'\|_2^2
 -\sum_{s\in R}h_s|u_-(s)|^2\leq0.
\]
在 \([\varepsilon,t]\) 上积分，得到 \(\|u_-(t)\|_2\leq\|u_-(\varepsilon)\|_2\)。负部映射在 \(L^2\) 中是 Lipschitz 的，而 \(u(\varepsilon)\to f\) 于 \(L^2\)，令 \(\varepsilon\downarrow0\) 得 \(\|u_-(t)\|_2\leq\|f_-\|_2\)。零初始负部于是始终为零。论证只使用一次空间弱导数，不要求负部在零点具有经典二阶导数。

\ProseParagraphBreak

全空间的纯热核解还保持总质量。若 \(f\in L^1\)，由绝对可积换序有 \(\int u(t)=\int f\)。对非负 \(f\)，当 \(0<s<t\) 时，半群关系及热核非负、总质量为一给出 \(\|u(t)\|_\infty\leq\|u(s)\|_\infty\)。降低的是全空间的峰值，不能据此说每一点的温度都随时间下降：初温支集外的零温点在正时间会升温。热量在空间中重新分布，而总量保持不变。若初值还有有限二阶绝对矩，Gauss 核的零一阶矩与每方向二阶矩 \(2\kappa t\) 给出
\[
 \int|x|^2u(t,x)\,\mathrm{d}x
 =\int|y|^2f(y)\,\mathrm{d}y+2d\kappa t\int f(y)\,\mathrm{d}y.
\]
当 \(f\geq0\) 时，这是总热量向半径 \(\sqrt t\) 的尺度扩散的定量表达。
